\documentclass[a4paper, serif, 10pt]{moderncv}

%% moderncv
% color and style
\moderncvcolor{black}
\moderncvstyle{banking}

%% page layout/margins
\usepackage[top=2.5cm, bottom=2.5cm, left=1.5cm, right=1.5cm]{geometry}

\nopagenumbers{}

%% Babel config for Spanish language
% \usepackage[spanish]{babel} has some conflicting issues with moderncv
% so we have to use enable the following options to make it work
%
% ref:
% - https://tex.stackexchange.com/a/140161/36007
\usepackage[spanish,es-lcroman]{babel}

%% fontspec
\usepackage{fontspec}

\IfFontExistsTF{Linux Libertine}{
  \setmainfont[Ligatures={TeX, Common}, Numbers=OldStyle]{Linux Libertine}
}{}
\IfFontExistsTF{Linux Libertine O}{
  \setmainfont[Ligatures={TeX, Common}, Numbers=OldStyle]{Linux Libertine O}
}{}

%% CTeX
% CJK support, used to show CJK characters in the resume
%
% - heading=false: disable ctex heading
% - punct=kaiming: use kaiming punctuations styles for CJK
% - scheme=plain: use plain scheme, do not override \normalsize font size
% - space=auto: space settings for CJK characters
%
% CJK fonts are resolved per locale language, e.g. Noto Serif/Sans CJK JP
% for Japanese resumes, with fallback to the next region if not installed
%
% ref:
% - http://ctan.mirrorcatalogs.com/language/chinese/ctex/ctex.pdf
\usepackage[UTF8, heading=false, punct=kaiming, scheme=plain, space=auto]{ctex}

\IfFontExistsTF{Noto Serif CJK SC}{
  \setCJKmainfont{Noto Serif CJK SC}
}{}
\IfFontExistsTF{Noto Sans CJK SC}{
  \setCJKsansfont{Noto Sans CJK SC}
}{}

%% line spacing
\usepackage{setspace}
\setstretch{1.125}

%% URL styling - use normal text instead of monospace
\urlstyle{same}

%% Patch \httplink and \httpslink to handle full URLs
% The original moderncv macros always prepend http:// or https:// to URLs.
% This causes issues when the URL already has a protocol (e.g., https://example.com)
% resulting in malformed URLs like https://https://example.com.
% This patch checks if the URL already contains "://" and uses it directly if so.
% Note: We use \str_set:Nx (x-expansion) to expand macros like \@homepage first.
\ExplSyntaxOn
\str_new:N \g_yamlresume_url_str

\RenewDocumentCommand{\httpslink}{O{}m}{
  \str_set:Nx \g_yamlresume_url_str {#2}
  \str_if_in:NnTF \g_yamlresume_url_str {://}
    { \url{#2} }
    { \url{https://#2} }
}

\RenewDocumentCommand{\httplink}{O{}m}{
  \str_set:Nx \g_yamlresume_url_str {#2}
  \str_if_in:NnTF \g_yamlresume_url_str {://}
    { \url{#2} }
    { \url{http://#2} }
}
\ExplSyntaxOff

\name{Lucía Fernández Morales}{}
\title{Enfermera Registrada | Urgencias y Cuidados Críticos}
\phone[mobile]{+34 612 345 678}
\email{lucia.fernandez@enfermeria.es}
\homepage{https://www.luciafernandez-enfermeria.es}

\address{Calle de la Salud 45, 3.º B, Madrid, Comunidad de Madrid, España, 28012}{}{}

\extrainfo{{\small \faLinkedin}\ \href{https://www.linkedin.com/in/lucia-fernandez-enfermeria}{@lucia-fernandez-enfermeria} {} {} {} • {} {} {} 
{\small \faTwitter}\ \href{https://twitter.com/lucia\_enfermera}{@lucia\_enfermera}}

\begin{document}

\maketitle

\section{Información básica}

\cvline{}{Enfermera registrada con más de siete años de experiencia en hospitales de alta complejidad, especializada en urgencias, cuidados críticos y atención centrada en la persona.
%
\begin{itemize}
\item Valoración integral del paciente y ejecución de cuidados de enfermería basados en evidencia.
\item Experiencia en triaje, monitorización hemodinámica, soporte vital avanzado y administración segura de medicamentos.
\item Coordinación con equipos multidisciplinares y comunicación clara con pacientes y familias.
\item Compromiso con la seguridad del paciente, la formación continua y la mejora de protocolos asistenciales.
\end{itemize}}
\section{Educación}

\cventry{sept 2014–jun 2018}
        {Licenciatura, Enfermería, Puntuación: 8,7/10}
        {Universidad Complutense de Madrid}
        {\url{https://www.ucm.es}}
        {}
        {\begin{itemize}
\item Formación generalista con prácticas clínicas en hospitalización, urgencias y atención primaria.
\item Trabajo de fin de grado sobre la prevención de infecciones asociadas a la atención sanitaria.
\item Participación en simulacros de soporte vital y en programas de educación para la salud.
\end{itemize}
\textbf{Cursos}: Fundamentos de Enfermería, Fisiología Humana, Farmacología Clínica, Cuidados Intensivos, Salud Pública, Ética y Legislación Sanitaria, Enfermería Materno-Infantil, Urgencias y Emergencias}
\section{Experiencia laboral}

\cventry{ene 2021 hasta la fecha}
        {Enfermera de Urgencias}
        {Hospital Universitario La Paz}
        {\url{https://www.comunidad.madrid/hospital/lapaz}}
        {}
        {\begin{itemize}
\item Realización de triaje inicial y clasificación de pacientes según prioridad asistencial en un servicio con más de 300 urgencias diarias.
\item Ejecución de cuidados de enfermería en pacientes críticos: monitorización, canalización de vías, administración de fármacos y apoyo en soporte vital avanzado.
\item Coordinación con médicos, técnicos y auxiliares para garantizar una atención rápida, segura y humanizada.
\item Registro riguroso de datos clínicos y educación al paciente y a la familia sobre el plan de cuidados y los signos de alarma.
\end{itemize}
\textbf{Palabras clave}: Triaje, Urgencias, RCP avanzada, Monitorización, Registro clínico}

\cventry{sept 2018–dic 2020}
        {Enfermera de Hospitalización}
        {Hospital Clínico San Carlos}
        {\url{https://www.comunidad.madrid/hospital/clinicosancarlos}}
        {}
        {\begin{itemize}
\item Atención integral a pacientes adultos en unidades de medicina interna y cirugía, con una media de 12 pacientes por turno.
\item Administración segura de medicamentos, curas, sondajes y control de signos vitales conforme a los protocolos del centro.
\item Colaboración en la preparación del paciente para pruebas diagnósticas e intervenciones quirúrgicas.
\item Educación sanitaria a pacientes y cuidadores para favorecer la continuidad de cuidados y la adherencia al tratamiento.
\end{itemize}
\textbf{Palabras clave}: Cuidados postoperatorios, Administración de medicamentos, Educación al paciente, Seguridad del paciente, Trabajo en equipo}
\section{Idiomas}

\cvline{Español}{Competencia nativa o bilingüe}
\cvline{Inglés}{Competencia profesional plena \hfill \textbf{Palabras clave}: B2-C1, Inglés sanitario}
\cvline{Francés}{Competencia limitada de trabajo \hfill \textbf{Palabras clave}: A2}
\section{Competencias}

\cvline{Cuidados de Enfermería}{Experto \hfill \textbf{Palabras clave}: Valoración del paciente, Cuidados básicos, Curas y vendajes, Sondajes, Control de signos vitales}
\cvline{Urgencias y Emergencias}{Avanzado \hfill \textbf{Palabras clave}: Triaje Manchester, RCP avanzada, Manejo de vías, Monitorización}
\cvline{Gestión Clínica}{Intermedio \hfill \textbf{Palabras clave}: Historia clínica electrónica, Registro de datos, Coordinación de equipos, Mejora de protocolos}
\cvline{Educación Sanitaria}{Avanzado \hfill \textbf{Palabras clave}: Educación al paciente, Promoción de la salud, Adherencia al tratamiento, Comunicación clínica}
\section{Premios}

\cventry{dic 2023}
        {Hospital Universitario La Paz}
        {Premio a la Excelencia Asistencial}
        {}
        {}
        {Reconocimiento por la mejora del protocolo de comunicación en el triaje de urgencias y por la formación de nuevos profesionales.}

\cventry{jun 2018}
        {Universidad Complutense de Madrid}
        {Mejor Expediente de Enfermería}
        {}
        {}
        {Distinción otorgada al expediente académico con mayor calificación de la promoción en el Grado en Enfermería.}
\section{Certificados}

\cventry{mar 2022}
        {Sociedad Española de Medicina de Urgencias y Emergencias}
        {Soporte Vital Avanzado (SVA)}
        {\url{https://www.semes.org}}
        {}
        {}

\cventry{may 2021}
        {Grupo Español de Triaje}
        {Curso de Triaje Manchester}
        {\url{https://www.triajemanchester.es}}
        {}
        {}

\cventry{oct 2020}
        {Hospital Universitario La Paz}
        {Enfermería de Urgencias y Emergencias}
        {\url{https://www.comunidad.madrid/hospital/lapaz}}
        {}
        {}
\section{Publicaciones}

\cventry{nov 2023}
        {Protocolo de actuación en triaje de urgencias pediátricas}
        {Revista Española de Enfermería de Urgencias}
        {\url{https://www.enfermeriaurgencias.es/articulos/triaje-pediatrico}}
        {}
        {\begin{itemize}
\item Artículo basado en la experiencia del servicio de urgencias del Hospital Universitario La Paz.
\item Propone criterios de valoración rápida y comunicación estructurada entre enfermería y pediatría.
\item Incluye indicadores de calidad para evaluar la adherencia al protocolo.
\end{itemize}}
\section{Referencias}

\cventry{\emaillink[marta.ruiz@salud.madrid.org]{marta.ruiz@salud.madrid.org}}
        {Supervisora de Urgencias}
        {Dra. Marta Ruiz Salgado}
        {+34 915 123 456}
        {}
        {Lucía destaca por su criterio clínico, su capacidad de priorización y su trato cercano con pacientes y familias. Es una profesional fiable y comprometida con la mejora continua.}
\section{Proyectos}

\cventry{feb 2022–nov 2022}
        {Programa comunitario para mejorar el autocuidado de personas con diabetes tipo 2.}
        {Programa de educación para la salud en diabetes}
        {\url{https://www.madrid.es/salud/diabetes}}
        {}
        {\begin{itemize}
\item Diseño de talleres sobre alimentación, ejercicio y control de la glucemia.
\item Coordinación con atención primaria y asociaciones de pacientes.
\item Evaluación de la satisfacción y de la adherencia a las recomendaciones.
\end{itemize}
\textbf{Palabras clave}: Educación sanitaria, Diabetes, Atención primaria, Autocuidado}
\section{Intereses}

\cvline{Voluntariado sanitario}{Cruz Roja, Primeros auxilios, Formación}
\cvline{Deporte}{Natación, Carrera, Yoga}
\cvline{Lectura}{Divulgación científica, Novela histórica}
\section{Voluntariado}

\cventry{jun 2016–sept 2019}
        {Voluntaria de Socorros y Emergencias}
        {Cruz Roja Española}
        {\url{https://www.cruzroja.es}}
        {}
        {\begin{itemize}
\item Participación en dispositivos preventivos y atención de primeros auxilios en eventos públicos.
\item Formación a nuevos voluntarios en soporte vital básico y manejo inicial de emergencias.
\item Colaboración en campañas de promoción de la salud y donación de sangre.
\end{itemize}}
    
\end{document}